\documentclass[t]{beamer}

\usepackage[utf8]{inputenc}
\usepackage[T1]{fontenc}
\usepackage{polski}
\usepackage{graphicx} 
\usepackage{epstopdf}
\usepackage{listings}
\usepackage[font={scriptsize,it}, labelfont={scriptsize,up,bf}]{caption}


\definecolor{mint}{rgb}{0.96, 1, 0.98}
\lstset
{
    language=[LaTeX]TeX,
    inputencoding=utf8,
    breaklines=true,
    basicstyle=\tt\scriptsize,
    keywordstyle=\color{red},
    identifierstyle=\color{blue},
    frame=lines,
    backgroundcolor=\color{mint}
}

%%%Komendy
%%%Beamer
\newcommand{\beamer}{
{\rmfamily\textsc{Beamer}}
}

\usetheme{Copenhagen}
\usecolortheme{crane}
\setbeamertemplate{navigation symbols}{} 
\expandafter\def\expandafter\insertshorttitle\expandafter{%
  \insertshorttitle\hfill%
  \insertframenumber\,/\,\inserttotalframenumber}
\mode<presentation>
\title{Próbna obrona pracy inżynierskiej}
\subtitle{Temat: Aplikacja WYSIWYG do tworzenia prezentacji multimedialnych z wykorzystaniem klasy dokumentów \beamer \LaTeX-a}
\author[J.A. Kumor]{Jan Adam Kumor}
\institute[IAiIS]{\begin{minipage}[c]{2.0cm}
			\includegraphics[width = 1.5cm]{IAiISLogo.eps}
		\end{minipage} Instytut Automatyki i Informatyki Stosowanej\\[0.6cm]}
\date[2015-10-23]{23 października 2015}

\usepackage{graphicx}

\begin{document}

\AtBeginSection{
	\begin{frame}
		\frametitle{Konspekt}
		\tableofcontents[currentsection, currentsubsection]
	\end{frame}
}

\begin{frame}
	\titlepage
\end{frame}

\begin{frame}[t]
	\frametitle{Temat pracy}
	\begin{description}[style=multiline]
		\item[Autor:] Jan Kumor, semestr 7. specjalność SID
		\item[Opiekun:] prof. nzw. dr hab. inż. Maciej Ławryńczuk \\ Zespół Technik Sterowania \\ Zakład Automatyki i Inżynierii Oprogramowania \\ Instytut Automatyki i Informatyki Stosowanej
		\item[Temat:] \emph{Aplikacja WYSIWYG do tworzenia prezentacji multimedialnych z wykorzystaniem klasy dokumentów} \beamer{} \LaTeX-a
	\end{description}
\end{frame}

\section{Cel pracy i możliwe zastosowania}

\begin{frame}
	\frametitle{Główne założenia}
	\begin{itemize}
		\item Graficzne środowisko do tworzenia prezentacji multimedialnych
		\item Zaadaptowanie systemu \LaTeX\ i klasy dokumentów \beamer\
		\item Systemy WYSIWYG (\emph{What You See Is What You Get})
	\end{itemize}
\end{frame}

\begin{frame}
	\frametitle{Możliwe zastosowania}
	\begin{itemize}
		\item Tworzenie dokumentów klasy \beamer\ bez znajomości \LaTeX-a
		\item Szybkie prototypowanie prezentacji
		\item Nauka \LaTeX\
	\end{itemize}
\end{frame}

\section{Istniejące rozwiązania}

\begin{frame}[T]
	\frametitle{Przykład 1. Gummi \footnote{\url{http://gummi.midnightcoding.org/}}}
		\begin{columns}[T]
			\begin{column}{.15\textwidth}
				\includegraphics[width=2cm]{gummi_logo.png}
			\end{column}
			\begin{column}{.85\textwidth}
				\begin{itemize}
					\item Typowe dla dokumentów \LaTeX\ workflow - plik \emph{.tex} i podgląd wyniku kompilacji
					\item Automatyczna kompilacja dokumentu po każdej wprowadzonej zmianie
					\item Dostępne  podstawowe operacje na tekście:
					\begin{itemize}
						\item \textbf{pogrubienie}, \emph{kursywa} oraz \underline{podkreślenie} tekstu
						\item wyrównanie tekstu do lewej i prawej strony oraz wyśrodkowanie
					\end{itemize}
					\item Wsparcie GUI przy wstawianiu grafiki, tabel, macierzy i bibliografii
				\end{itemize}
			\end{column}		
		\end{columns}
\end{frame}

\begin{frame}[T]
	\frametitle{Przykład 1. Gummi}
	\begin{figure}[p]
   	    \centering
		\includegraphics[width=.95\textwidth]{gummi_prtscn.png}
   	    \caption{Okno programu Gummi}
	\end{figure}
\end{frame}

\begin{frame}[t]
	\frametitle{Przykład 2. BaKoMa \TeX\ (True WYSIWYG \LaTeX\ System)\footnote{\url{http://www.bakoma-tex.com/}}}
		\begin{columns}[T]
			\begin{column}{.15\textwidth}
				\includegraphics[width=2cm]{bakoma_logo.png}
			\end{column}
			\begin{column}{.85\textwidth}
				\begin{itemize}
					\item Pakiet aplikacji komercyjnych 
					\item Wsparcie GUI dla większości funckji \LaTeX\ 
					\item Obsługa wielu klas dokumentów z poziomu GUI
					\item Prawdziwy WYSIWYG w czasie rzeczywistym
					\item W tle działa silnik \TeX - tylko dla aktualnie zmienianej strony
				\end{itemize}
			\end{column}		
		\end{columns}
\end{frame}

\begin{frame}[t]
	\frametitle{Przykład 2. BaKoMa \TeX\ (True WYSIWYG \LaTeX\ System)}
		\begin{figure}[p]
     	    \centering
			\includegraphics[width=.95\textwidth]{bakoma_prtscn.png}
     	    \caption{Okno programu BaKoMa Word}
		\end{figure}
		
\end{frame}

\section{Zapropownowane rozwiązanie i wynik pracy}

\begin{frame}
	\frametitle{Aplikacja Rapidbeamer}
	\begin{itemize}
		\item język programowania - Python 3.5
		\item framework Kivy
	\end{itemize}
	\begin{figure}[htp]
		\centering
		\includegraphics[width=.80\textwidth]{rapidbeamer.png}
		\caption{Okno aplikacji Rapidbeamer}
		\label{rbm}
	\end{figure}
\end{frame}

\begin{frame}
	\frametitle{Architektura rozwiązania}
	Wzorzec MVC pozwala na wygodne rozdzielenie zakresów odpowiedzialności poszczególnych części aplikacji.\\
	\begin{center}
		\includegraphics[width=0.75\textwidth]{mvc.png}
	\end{center}
\end{frame}

\begin{frame}
	\frametitle{Warstwa M - Model}
	\begin{itemize}
		\item Python 3.5
		\item Trójwarstwowa strukrura modelu
			\begin{itemize}
				\alert<1>{\item Warstwa abstrakcyjna}
				\alert<2>{\item Warstwa dokumentu \TeX\ }
				\alert<3>{\item Warstwa generowanego kodu}
			\end{itemize}
	\end{itemize}
	\hspace*{-11mm}
	\includegraphics<1-3>[width=\paperwidth]{model_layers.png}
	\includegraphics<4>[width=\paperwidth]{model_layers_arrows.png}
	
\end{frame}

\begin{frame}
	\frametitle{Warstwa M - Model}
	Realizuje większość logiki aplikacji
	\begin{itemize}
		\item Generowanie kodu \LaTeX\ - wzorzec kompozytu
		\item Zapis i odczyt projektu - format JSON				
		\item API: metody manipulacji modelem + wzorzec obserwatora
	\end{itemize}
	\hspace*{-7mm}
	\begin{minipage}{\paperwidth}
		\begin{minipage}{.5\paperwidth}
			\begin{figure}
				\includegraphics[width=.4\paperwidth]{composite.png}
				\caption{"Composite UML class diagram"\\autorstwa Trashtoyderivative}
				\label{cap}
			\end{figure}
		\end{minipage}
		\begin{minipage}{.5\paperwidth}
			\begin{figure}
				\includegraphics[width=.4\paperwidth]{observer.png}
				\caption{"Observer classes pl"\\autorstwa Zyxist}
				\label{obs}
			\end{figure}
		\end{minipage}
	\end{minipage}
\end{frame}

\begin{frame}
	\frametitle{Warstwa VC - Widok i kontroler}
	Umożliwia interakcję użytkownika z aplikacją.
	\begin{itemize}
		\item Sterowanie zdarzeniami
		\item Dwa współpracujące wątki: widoku i kontrolera
		\item widok zaimplementowany z użyciem frameworku Kivy
	\end{itemize}
\end{frame}

\begin{frame}
	\frametitle{Warstwa VC - Widok i kontroler cd}

	\raisebox{-.4cm}{\includegraphics[width=1.2cm]{kivy_logo.png}} Framework Kivy\footnote{http://kivy.org/}

	\begin{itemize}
		\item Otwarta biblioteka języka Python do tworzenia graficznego interfejsu użytkownika 
		\item Wieloplatformowość: \emph{Linux}, \emph{Windows}, \emph{OS X}, \emph{Android} i \emph{iOS}
		\item Szerokie wsparcie dla NUI (\emph{Natural User Interfaces})
		\item Dedykowany język opisu interfejsu \emph{kivy language}
	\end{itemize}
\end{frame}

\begin{frame}[plain]
	\frametitle{Budowa interfejsu}
	\hspace*{-11mm}
	\includegraphics[width=\paperwidth]{rapidbeamer.png}
\end{frame}

\begin{frame}
	\frametitle{Moduł prezentacji}
	\centering \alert{Stan zaawansowania pracy nie pozwala powiedzieć aktualnie nic na ten temat...}
\end{frame}

\section{Wnioski i podsumowanie}

\begin{frame}
	\frametitle{Wnioski i podsumowanie}
	\begin{itemize}
		\item Udało się stworzyć aplikację umożliwiającą wygodne i szybkie tworzenie prezentacji multimedialnych.
		\item \alert{Więcej wniosków mogłoby dotyczyć testów wydajnościowych generatora kodu lub mechanizmu prezentacji}
	\end{itemize}
\end{frame}

\begin{frame}[c]
	\Huge Dziękuję
\end{frame}

\end{document}
